% Suggested final budget: 0.75--1.0 page.
\section{Experimental Design}
\label{sec:experiments}
\subsection{Data-size study}
\todo{Confirm completed $N\in\{25,50,100,200\}$ cells and any justified $N=400$
extension. State training seed counts, nested-subset construction, fixed
optimizer-step budget, and shared test initial conditions. Distinguish planned
experiments from completed experiments. Explain the repeated exposure implied
by a fixed-step rather than fixed-epoch budget.}

\subsection{Backbone study}
\todo{List the task subset, demonstration count per task, completed architectures,
seed counts, and settings held constant. State parameter counts and any tuning
or compute differences. If capacities differ, restrict conclusions to these
specific implementations rather than claiming a pure architecture effect.}

\subsection{Evaluation and statistical reporting}
\todo{Confirm actual evaluation protocol: held-out seeds, number of episodes,
simulation backend, inference randomness, parallel environments, episode
termination behavior, and checkpoint selection. Keep validation, training-seed
diagnostics, and held-out test results separate.}

For a single training run with $E$ evaluation episodes, report the success
count $S$ and rate $\hat p=S/E$. Define \metriconce{} as success at any point
in an episode and \metricend{} as success at its final step. Across $R$
training seeds, summarize the seed-level rates with
\begin{equation}
  \bar p=\frac{1}{R}\sum_{r=1}^{R}\hat p_r,
  \qquad
  s_p=\sqrt{\frac{1}{R-1}\sum_{r=1}^{R}(\hat p_r-\bar p)^2}.
  \label{eq:seed-statistics}
\end{equation}
\todo{State which interval/comparison procedure was actually used, and its
resampling unit. Seed-level standard deviation is not a confidence interval.
For a single training seed, do not report a fabricated standard deviation.
Document the hardware and timing method for efficiency measurements.}
